\documentclass[aspectratio=169,10pt]{beamer}

\usetheme{Madrid}
\usecolortheme{default}

\usepackage{booktabs}
\usepackage{array}

\definecolor{pgoBlue}{RGB}{37,84,150}
\definecolor{pgoGreen}{RGB}{54,130,91}
\definecolor{pgoOrange}{RGB}{191,104,43}

\setbeamercolor{structure}{fg=pgoBlue}
\setbeamertemplate{navigation symbols}{}
\setbeamertemplate{caption}[numbered]

\newcommand{\case}[1]{\texttt{\detokenize{#1}}}
\newcommand{\pct}[1]{#1\%}
\newcommand{\good}[1]{\textcolor{pgoGreen}{#1}}
\newcommand{\warn}[1]{\textcolor{pgoOrange}{#1}}

\title[Bunny Dynamic FEM]{Bunny Dynamic Drop on a Conservative Cubic Domain}
\subtitle{Cubic-Linear vs. Refined Cubic-Linear vs. Tricubic Hermite}
\author{libpgo experiments}
\date{\today}

\begin{document}

\begin{frame}
  \titlepage
\end{frame}

\begin{frame}{Question}
  \begin{block}{Main question}
    During a contact-rich dynamic drop, does tricubic Hermite produce a visible-surface trajectory closer to a same-domain tet-linear numerical reference than cubic-linear FEM?
  \end{block}

  \vspace{0.5em}
  \begin{itemize}
    \item Compare original cubic-linear, refined cubic-linear, and cubic-tricubic-Hermite.
    \item Use a same-domain tet-linear solve as the numerical reference.
    \item Include \case{cubic_linear_x8} as a comparable-DOF low-order baseline.
    \item Report accuracy and contact timing only.
  \end{itemize}
\end{frame}

\begin{frame}{Cases}
  \small
  \begin{table}
    \centering
    \begin{tabular}{llp{0.48\linewidth}}
      \toprule
      \normalfont case & formulation & role \\
      \midrule
      \case{tet_ref} & tet-linear & same-domain numerical reference \\
      \case{cubic_linear} & cubic-linear & original r15 low-order cubic baseline \\
      \case{cubic_linear_x8} & cubic-linear & same domain; each cube split 2x2x2 \\
      \case{cubic_hermite} & cubic-tricubic-Hermite & high-order method under test \\
      \bottomrule
    \end{tabular}
  \end{table}

  \vspace{0.5em}
  \begin{block}{Why include \case{cubic_linear_x8}?}
    It checks whether the Hermite gain is only a DOF-count effect.
  \end{block}
\end{frame}

\begin{frame}{Setup}
  \small
  \begin{columns}[T]
    \begin{column}{0.48\linewidth}
      \textbf{Dynamic setup}
      \begin{itemize}
        \item Stable Neo-Hookean material
        \item Gravity: $(0, -9.81, 0)$
        \item Initial velocity: $(0,0,0)$
        \item Timestep: $0.001$
        \item Steps: $800$
        \item Dump interval: $10$
        \item IPC contact with ground obstacle
      \end{itemize}
    \end{column}
    \begin{column}{0.48\linewidth}
      \textbf{DOF scale}
      \begin{table}
        \centering
        \begin{tabular}{lr}
          \toprule
          case & DOFs \\
          \midrule
          \case{tet_ref} & 363,714 \\
          \case{cubic_linear} & 9,102 \\
          \case{cubic_linear_x8} & 62,355 \\
          \case{cubic_hermite} & 72,816 \\
          \bottomrule
        \end{tabular}
      \end{table}
    \end{column}
  \end{columns}
\end{frame}

\begin{frame}{Surface Trajectory Metric}
  \small
  Let $S$ be the visible surface vertices, $x_i^0$ the rest position,
  $x_i^c(t)$ the position for case $c$, and $x_i^{ref}(t)$ the tet-reference
  position at dumped frame $t$.

  \[
    u_i^c(t) = x_i^c(t) - x_i^0,
    \qquad
    u_i^{ref}(t) = x_i^{ref}(t) - x_i^0
  \]

  \[
    E_c(t) =
    \frac{
      \sqrt{\sum_{i \in S}\|u_i^c(t)-u_i^{ref}(t)\|_2^2}
    }{
      \sqrt{\sum_{i \in S}\|u_i^{ref}(t)\|_2^2}
    }
  \]

  \vspace{0.5em}
  \begin{block}{Interpretation}
    $E_c(t)$ measures surface displacement trajectory error relative to the
    same-domain tet reference. It is computed on the original visible bunny
    surface, not on volume mesh nodes.
  \end{block}
\end{frame}

\begin{frame}{Windowed Metrics}
  \small
  \begin{table}
    \centering
    \begin{tabular}{lp{0.62\linewidth}}
      \toprule
      window & dumped frames \\
      \midrule
      full & $10,20,\ldots,790$ \\
      pre-contact & $100,110,\ldots,390$ \\
      post-contact & $410,420,\ldots,790$ \\
      late & $600,610,\ldots,790$ \\
      \bottomrule
    \end{tabular}
  \end{table}

  \vspace{0.8em}
  \begin{itemize}
    \item \case{tet_ref} reaches the impact threshold at frame 410.
    \item The late window isolates accumulated post-impact trajectory error.
    \item Frame 0 is not used as a headline relative-error metric because the
      reference displacement norm is very small.
  \end{itemize}
\end{frame}

\begin{frame}{Trajectory Accuracy}
  \small
  \begin{table}
    \centering
    \begin{tabular}{llrrr}
      \toprule
      window & statistic & cubic-linear & x8 linear & Hermite \\
      \midrule
      full & mean rel L2 & \pct{3.30} & \pct{1.34} & \good{\pct{0.63}} \\
      full & p95 rel L2 & \pct{10.32} & \pct{2.92} & \good{\pct{1.29}} \\
      pre & mean rel L2 & \pct{0.19} & \pct{0.13} & \good{\pct{0.12}} \\
      pre & p95 rel L2 & \pct{0.46} & \pct{0.28} & \good{\pct{0.23}} \\
      post & mean rel L2 & \pct{4.25} & \pct{1.42} & \good{\pct{0.35}} \\
      post & p95 rel L2 & \pct{10.32} & \pct{2.74} & \good{\pct{0.93}} \\
      late & mean rel L2 & \pct{6.33} & \pct{2.06} & \good{\pct{0.53}} \\
      late & p95 rel L2 & \pct{11.47} & \pct{2.92} & \good{\pct{1.10}} \\
      \bottomrule
    \end{tabular}
  \end{table}

  \vspace{0.8em}
  \begin{itemize}
    \item Full window: dumped frames $10,20,\ldots,790$.
    \item Refinement removes much of the original cubic-linear error.
    \item Hermite remains closest to the tet reference at comparable DOF scale.
  \end{itemize}
\end{frame}

\begin{frame}{Before vs. After Contact}
  \small
  \begin{table}
    \centering
    \begin{tabular}{lrrr}
      \toprule
      case & pre-contact & post-contact & post/pre \\
      \midrule
      \case{cubic_linear} & \pct{0.19} & \pct{4.25} & 22.8x \\
      \case{cubic_linear_x8} & \pct{0.13} & \pct{1.42} & 10.8x \\
      \case{cubic_hermite} & \good{\pct{0.12}} & \good{\pct{0.35}} & \good{2.9x} \\
      \bottomrule
    \end{tabular}
  \end{table}

  \vspace{0.8em}
  \begin{block}{Main observation}
    The experiment is primarily testing contact response. Before impact, all
    cubic methods are close to the tet reference; after impact, the low-order
    trajectories separate more strongly.
  \end{block}
\end{frame}

\begin{frame}{Contact Timing}
  \small
  \begin{table}
    \centering
    \begin{tabular}{lr}
      \toprule
      case & impact frame \\
      \midrule
      \case{tet_ref} & 410 \\
      \case{cubic_linear} & \warn{400} \\
      \case{cubic_linear_x8} & 410 \\
      \case{cubic_hermite} & 410 \\
      \bottomrule
    \end{tabular}
  \end{table}

  \vspace{0.8em}
  \begin{itemize}
    \item The original cubic-linear case reaches the contact threshold one
      dump interval earlier than the tet reference.
    \item The refined cubic-linear and Hermite cases match the tet reference
      impact frame.
  \end{itemize}
\end{frame}

\begin{frame}{Error Reduction}
  \small
  \begin{table}
    \centering
    \begin{tabular}{lr}
      \toprule
      comparison & reduction \\
      \midrule
      x8 late error vs original cubic-linear & 3.1x lower \\
      Hermite late error vs x8 cubic-linear & \good{3.9x lower} \\
      \bottomrule
    \end{tabular}
  \end{table}

  \vspace{0.8em}
  \begin{itemize}
    \item \case{cubic_linear_x8} is a strong baseline.
    \item Hermite still improves over it on post-contact trajectory accuracy.
  \end{itemize}
\end{frame}

\begin{frame}{Limitations}
  \small
  \begin{itemize}
    \item The tet reference is a same-domain numerical baseline, not ground truth.
    \item Metrics are evaluated only at dumped frames every 10 simulation steps.
    \item \case{cubic_linear_x8} is comparable-DOF, not exactly equal-DOF.
    \item The late window is diagnostic: it highlights accumulated post-impact error.
  \end{itemize}
\end{frame}

\begin{frame}{Conclusion}
  \small
  \begin{block}{Main result}
    In the bunny dynamic drop experiment, tricubic Hermite is closer to the
    same-domain tet-linear reference than both original cubic-linear and
    2x2x2 refined cubic-linear, especially after contact.
  \end{block}

  \vspace{0.7em}
  \begin{itemize}
    \item Post-contact mean rel L2: \pct{4.25} $\rightarrow$ \pct{1.42} $\rightarrow$ \good{\pct{0.35}}.
    \item Late mean rel L2: \pct{6.33} $\rightarrow$ \pct{2.06} $\rightarrow$ \good{\pct{0.53}}.
    \item The advantage is trajectory accuracy under contact.
  \end{itemize}
\end{frame}

\begin{frame}{Reproducibility Pointers}
  \scriptsize
  \begin{itemize}
    \item Report: \texttt{examples/scripts/tricubic-hermit-dynamic-compare/report/bunny\_dynamic\_drop\_compare\_conservative\_r15.md}
    \item Runner: \texttt{examples/scripts/tricubic-hermit-dynamic-compare/dynamic\_compare.py}
    \item Output root: \texttt{examples/outputs/bunny-dynamic-drop-compare-conservative-r15}
    \item Metrics use dumped surfaces: \texttt{<case>/surface/surface0000.obj} through \texttt{surface0790.obj}
  \end{itemize}
\end{frame}

\end{document}
